\subsection{相对连通映射}

定义 1. --- 设 X 与 Y 是拓扑空间，f 是从 X 到 Y 的连续映射。称映射 f 为相对连通的，如果 Y 的每一点都有一个由满足下列条件的集合 V 组成的基本邻域系：\(f^{-1}(V)\) 连通且非空。

设 \(f \colon X \to Y\) 是连续映射；为使 \(f\) 相对连通，必要且充分的是 \(Y\) 的每一点都有一个邻域 \(V\)，使映射 \(f_V \colon f^{-1}(V) \to V\)（由 \(f\) 导出）相对连通。

设 X 与 Y 是拓扑空间，\(f: X \to Y\) 是相对连通连续映射。f 的像在 Y 中稠密。对 Y 的每个开子集 V，映射 \(f_{\mathrm{V}}: f^{-1}(\mathrm{V}) \to \mathrm{V}\) 相对连通。

命题 4. --- 设 X 与 Y 是拓扑空间，\(f: X \to Y\) 是相对连通连续映射。

\begin{enumerate}
\def\labelenumi{\alph{enumi})}
\item
  对 X 的每个连通分支 U，\(f(\overline{\mathrm{U}})\) 是 Y 的连通、开且闭的子集，且我们有 \(f^{-1}(\overline{f(\mathrm{U})}) = \mathrm{U}\)。
\item
  对 Y 的每个连通分支 V，存在 X 的一个连通分支 U，使 \(V = \overline{f(U)}\)。由 f 取子空间导出的、从 U 到 V 的映射相对连通。
\item
  X（相应地，Y）的连通分支都是开且闭的。
\item
设 V 是由满足下列条件的点 \(y \in Y\) 组成的集合：它们有一个邻域 W，使 \(f^{-1}(W)\) 连通且与 U 相交；V 是 Y 的开子集。它包含 \(\overline{f(U)}\)，因为 f 相对连通。反之，设 \(y \in V\)；设 W 是 \(y\) 的一个邻域，使 \(f^{-1}(W)\) 连通且与 U 相交。对每个邻域 \(W'\)（\(y\) 的、使 \(f^{-1}(W')\) 连通的邻域），\(f^{-1}(W \cap W')\) 非空，因此 \(f^{-1}(W \cup W')\) 连通（TG, I, p. 81, 命题 2）。依假设，\(f^{-1}(W \cup W')\) 与 U 相交；于是 \(f^{-1}(W \cup W') \subset U\)，且更有 \(W' \cap f(U) \neq \varnothing\)。这证明了 \(y \in \overline{f(U)}\)；因此 \(V = \overline{f(U)}\)。特别地，集合 \(\overline{f(U)}\) 在 Y 中既开且闭；此外，命题 1（TG, I, p. 81）蕴含它连通。
\end{enumerate}

前面的论证还表明 \(f^{-1}(\overline{f(\mathrm{U})}) \subset \mathrm{U}\)，从而有等式 \(f^{-1}(\overline{f(\mathrm{U})}) = \mathrm{U}\)，另一个包含关系是显然的。特别地，U 在 X 中既开且闭。

\begin{enumerate}
\def\labelenumi{\alph{enumi})}
\setcounter{enumi}{1}
\tightlist
\item
设 V 是 Y 的一点 y 的连通分支，W 是 y 的一个使 \(f^{-1}(W)\) 连通非空的邻域，U 是包含 \(f^{-1}(W)\) 的 X 的连通分支。由于 \(y \in \overline{f(U)}\)，由 a) 及 y 的连通分支的定义得出 \(\mathrm{V} = \overline{f(\mathrm{U})}\)。因此 V 在 Y 中既开且闭。
\end{enumerate}

推论 1. --- 设 X 与 Y 是拓扑空间，\(f: X \to Y\) 是连续且相对连通的映射。通过取连通分支，映射 \(f\) 诱导一个从 X 的连通分支所成之集到 Y 的连通分支所成之集上的双射。

推论 2. --- 设 X 与 Y 是拓扑空间，\(f: X \to Y\) 是连续映射。为使 \(f\) 相对连通，必要且充分的是下列三条性质成立：

\begin{enumerate}
\def\labelenumi{\alph{enumi})}
\item
  像 f(X) 在 Y 中稠密；
\item
  空间 Y 局部连通；
\item
  对 Y 的每个开连通集合 V，集合 \(f(\mathrm{V})\) 连通。
\end{enumerate}

假设映射 f 相对连通。\(f(\mathrm{X})\) 在 Y 中的稠密性由相对连通映射的定义得出。设 V 是 Y 的开子集；映射 \(f_{\mathrm{V}}: f^{-1}(\mathrm{V}) \to \mathrm{V}\) 相对连通。由命题 4，V 的连通分支在 V 中既开且闭。由此得出 Y 局部连通（TG, I, p. 85, 命题 11）。为证明断言 c)，只需证明：若 Y 连通，则 X 连通。由引理，存在 X 的一个连通分支 U，使 \(Y = \overline{f(U)}\) 且 \(f^{-1}(\overline{f(U)}) = U\)，于是 U = X，故 X 连通。

反之，假设条件 \(a\)、\(b\)、\(c\) 满足，我们来证明 f 相对连通。设 \(y\) 是 Y 的一点；由于 Y 局部连通，\(y\) 有一个由连通开邻域组成的基本邻域系。若 W 是这样一个邻域，则条件 \(c\) 与 \(a\) 蕴含 \(f^{-1}(W)\) 连通且非空。因此映射 f 相对连通。

推论 3. --- 设 X 与 Y 是拓扑空间，\(f\colon X\to Y\) 是连续且相对连通的映射。设 F 是一个集合，\(g\colon X\to F\) 是局部常值映射。存在唯一的局部常值映射 \(h\colon Y\to F\)，使 \(g=h\circ f\)。

g 在 X 的每个连通分支上的限制是常值的。由推论 1，存在一个映射 \(h: Y \to F\)，它在 Y 的每个连通分支上取常值，使 \(g = h \circ f\)。映射 h 是局部常值的，因为 Y 的连通分支都是开的。这样一个映射的唯一性由 \(f(X)\) 在 Y 中稠密这一事实得出。

例. --- 1) 设 X 是拓扑空间，R 是 X 上的一个等价关系。记 Y 为商拓扑空间 X/R，记 \(f\colon X \to Y\) 为典范映射。假设 R 的等价类都连通。那么，对 Y 的每个开连通子集 V，集合 \(f^{-1}(V)\) 连通（TG, I, p. 23, 推论 1 及 p. 82, 命题 7）。若空间 Y 局部连通，则映射 f 因而相对连通。

\begin{enumerate}
\def\labelenumi{\arabic{enumi})}
\setcounter{enumi}{1}
\tightlist
\item
  设 Y 是局部道路连通拓扑空间，X 是 Y 的开子集。X 中道路的空间 \(\mathcal{C}_{c}(\mathbf{I};\mathbf{X})\) 与 Y 中道路的空间 \(\mathcal{C}_{c}(\mathbf{I};\mathbf{Y})\) 的一个子空间等同起来；假设它在其中稠密。于是 X 到 Y 中的典范单射是相对连通映射。
\end{enumerate}

事实上，只需证明：对 Y 的每个非空连通开子集 V，集合 \(V \cap X\) 连通且非空。空间 \(\mathcal{C}_{c}(\mathbf{I};\mathrm{V})\) 是 \(\mathcal{C}_{c}(\mathbf{I};\mathrm{Y})\) 的一个非空开集；它与 \(\mathcal{C}_{c}(\mathbf{I};\mathrm{X})\) 相交，这就证明了 \(V \cap X \neq \varnothing\)。设 x 与 \(x'\) 是 \(V \cap X\) 的点。由于 \(V \cap X\) 局部道路连通，点 x 与 \(x'\) 有开邻域 U 与 \(U'\)，它们道路连通且包含于 \(V \cap X\)。由于 V 道路连通（III, p. 260, 命题 7），存在 V 中连接 x 与 \(x'\) 的道路。由稠密性假设及紧收敛拓扑的定义，存在 \(V \cap X\) 中连接 U 的一点与 \(U'\) 的一点的一条道路。于是存在一条连接 x 与 \(x'\) 的、\(V \cap X\) 中的道路，这就证明了集合 \(V \cap X\) 连通。

命题 5. --- 设 X 与 Y 是拓扑空间，\(f: X \to Y\) 是连续且相对连通的映射。对 Y 的覆叠的每个偶 \((T, T')\)，映射 \(f^{*}: \mathcal{C}_{Y}(T; T') \to \mathcal{C}_{X}(X \times_{Y} T; X \times_{Y} T')\) 是双射的。

设 \(\mathcal{F}\) 是 X 上由从 \(X \times_{Y} T\) 到 \(X \times_{Y} T'\) 中的 X-态射构成的层，\(\mathcal{G}\) 是 Y 上由从 T 到 \(T'\) 中的 Y-态射构成的层（I, p. 45, 例 4）。对 Y 的每个开集 U，置 \(\varphi_{\mathrm{U}} = (f_{\mathrm{U}})^{*}: \mathcal{G}(\mathrm{U}) \to \mathcal{F}(f^{-1}(\mathrm{U}))\)。诸映射 \(\varphi_{\mathrm{U}}\) 定义一个层的态射 \(\varphi: \mathcal{G} \to \varphi_{*}(\mathcal{F})\)，只需证明 \(\varphi\) 是层的一个同构。

由于 Y 局部连通（IV, p. 354, 推论 2），T 与 T' 在其上可平凡化的连通开集构成 Y 的拓扑的一个基。由 I, p. 55 的推论 2，只需证明：对这样一个开集 U，映射 \(\varphi_{\mathrm{U}}\) 是双射的，这使我们可以假设 Y 连通，且覆叠 T 与 T' 是平凡覆叠 \(Y \times F\) 与 \(Y \times F'\)，其中 F 与 F' 是赋有离散拓扑的集合。映射 \((x, (y, t)) \mapsto (x, t)\) 把 X-空间 \(X \times_{Y} (Y \times F)\) 与 \(X \times F\) 等同起来（相应地，把 X-空间 \(X \times_{Y} (Y \times F')\) 与 \(X \times F'\) 等同起来）。由于空间 X 连通（IV, p. 354, 推论 2），集合 \(\mathcal{C}_{\mathrm{Y}}(Y \times F; Y \times F')\) 与 \(\mathcal{C}_{\mathrm{X}}(X \times F; X \times F')\) 二者都与从 F 到 F' 中的映射所成的集合 \(\mathcal{F}(F;F')\) 等同，且映射 \(f^{*}\) 与 \(\mathcal{F}(F;F')\) 的恒等映射等同。证明到此完成。

推论 1. --- 设 X 与 Y 是拓扑空间，\(f: X \to Y\) 是相对连通连续映射。设 T 与 T' 是 Y 的覆叠。若 X 的覆叠 \(X \times_Y T\) 与 \(X \times_Y T'\) 同构，则覆叠 T 与 T' 同构。

事实上，设 \(h \colon X \times_{Y} T \to X \times_{Y} T'\) 与 \(h' \colon X \times_{Y} T' \to X \times_{Y} T\) 是互为逆的 X-同构。由命题 5，存在 Y-态射 \(g \colon \mathrm{T} \to \mathrm{T}'\) 与 \(g' \colon \mathrm{T}' \to \mathrm{T}\)，使 \(f^{*}(g) = h\) 且 \(f^{*}(g') = h'\)。于是有 \(f^{*}(g' \circ g) = f^{*}(g') \circ f^{*}(g) = \mathrm{Id}_{\mathrm{X} \times_{\mathrm{Y}} \mathrm{T}}\)，因此 \(g' \circ g = \mathrm{Id}_{\mathrm{T}}\)，因为映射 \(f^{*}\) 是单射的。类似地，\(g \circ g' = \mathrm{Id}_{\mathrm{T}'}\)。因此覆叠 T 与 \(\mathrm{T}'\) 同构。

推论 2. --- 设 X 与 Y 是拓扑空间，\(f: X \to Y\) 是连续且相对连通的映射。若空间 X 单连通，则空间 Y 亦然。

事实上，由推论 1，Y 的每个覆叠都可平凡化。

推论 3. --- 设 X 是局部道路连通拓扑空间，Y 是可解圈拓扑空间，\(f: X \to Y\) 是相对连通连续映射。对 X 的每一点 \(x\)，同态 \(\pi_1(f, x): \pi_1(X, x) \to \pi_1(Y, f(x))\) 是满射的。

由 IV, p. 353 的命题 4，我们可以假设空间 X 与 Y 都连通。于是推论由 IV, p. 352 的命题 5 与命题 2 得出。

命题 6. --- 设 X 与 Y 是拓扑空间，\(f: \mathrm{X} \to \mathrm{Y}\) 是连续且相对连通的映射。假设 Y 的每一点都有一个开邻域 V，其逆像 \(f^{-1}(V)\) 单连通。那么，对 X 的每个覆叠 Z，存在 Y 的一个覆叠 T，使 \(X \times_{Y} T\) X-同构于 Z。

设 \(\mathcal{U}\) 是 Y 中其逆像在 X 中单连通的开子集所成的集合。对每个 \(V \in \mathcal{U}\)，覆叠 \(f^{-1}(V) \times_{X} Z\)（\(f^{-1}(V)\) 的覆叠）是可平凡化的；于是存在一个离散空间 \(F_V\) 及覆叠的一个同构 \(g_V: f^{-1}(V) \times_{X} Z \to f^{-1}(V) \times F_V\)。对开集的每个偶 \((V, V')\)（属于 \(\mathcal{U}\)），映射 \(f_{V \cap V'}: f^{-1}(V \cap V') \to V \cap V'\) 相对连通。由 IV, p. 356 的命题 5，存在 \(V \cap V'\) 的覆叠的唯一同构 \(h_{V',V}: (V \cap V') \times F_V \to (V \cap V') \times F_{V'}\)，使得有 \(f^*(h_{V',V})(x,t) = g_{V'}(g_V^{-1}(x,t))\)（对所有 \(x \in V \cap V'\) 及每个 \(t \in F_V\)）。若 V、\(V'\)、\(V''\) 是 \(\mathcal{U}\) 的元素，则有 \(h_{V'',V}(x,t) = h_{V'',V'}(h_{V',V}(x,t))\)（对所有 \(x \in V \cap V' \cap V''\) 及每个 \(t \in F_V\)）。于是存在唯一的 Y-空间 T，且对每个 \(V \in \mathcal{U}\)，存在同构 \(h_V: T_V \to V \times F_V\)，使得有 \(h_{V',V}(x,t) = h_{V'} \circ h_V^{-1}(x,t)\)——其中偶 \((V, V')\) 属于 \(\mathcal{U}\)、\(x \in V \cap V'\) 任意、\(t \in F_V\) 任意（参见 TG, I, p. 16）。

空间 T 特别地是 Y 的一个覆叠。此外，存在唯一的、到 Z 上的 \(X \times_{Y} T\) 的映射，它在 \(f^{-1}(V) \times_{Y} T\) 上的限制由 \(g_{V}^{-1} \circ f^{*}(h_{V})\) 给出，且它是 X-空间的一个同构，命题由此得证。

推论. --- 设 X 与 Y 是可解圈拓扑空间，\(f: X \to Y\) 是连续且相对连通的映射。假设 Y 的每一点都有一个邻域 V，其逆像 \(f^{-1}(V)\) 单连通。那么，对 X 的每一点 x，同态 \(\pi_{1}(f,x): \pi_{1}(X,x) \to \pi_{1}(Y,f(x))\) 是双射的。

可以假设空间 X 与 Y 都连通（IV, p. 353, 命题 4）。由推论 3（IV, p. 357），同态 \(\pi_1(f,x)\) 是满射的。命题 6 与 IV, p. 351 的命题 1 蕴含它是单射的。

注. --- 设 Y 是一个拓扑空间，X、X' 与 Y' 是 Y 的子空间，满足 X ⊂ X' ⊂ Y' ⊂ Y。假设 X 到 Y 中的典范单射相对连通。对 Y 的每个连通开子集 V，集合 V ∩ X 连通且在 V 中稠密（IV, p. 354, 推论 2）；于是集合 V ∩ X' 连通（TG, I, p. 81, 命题 1）。这证明了 X' 到 Y' 中的典范单射相对连通。

设 V 是 Y 的一个开子集；由前文，\(V \cap X\) 到 \(V \cap X'\) 中的典范单射相对连通。若集合 \(V \cap X\) 单连通，则 \(V \cap X'\) 亦然，这由 IV, p. 357 的推论 2 得出。因此，若 X 到 Y 中的典范单射满足命题 6 的假设，则 \(X'\) 到 \(Y'\) 中的典范单射亦然。

例. --- 设 Y 是一个局部有限维微分流形，Z 是 Y 的一个闭子流形（VAR, R, 5.8.3）。置 X = Y - Z。

\begin{enumerate}
\def\labelenumi{\alph{enumi})}
\item
  若 Z 在每点的余维数至少为 2，则 X 到 Y 中的典范单射相对连通。事实上，设 z 是 Z 的一点；存在 z 在 Y 中的开邻域 V 及 V 到 R 上的一个有限维向量空间 E 上的同胚 \(\varphi\)，使 \(\varphi(V \cap Z)\) 是 E 的一个余维数 \(\geqslant 2\) 的向量子空间 F。集合 E-F 连通（TG, VI, p. 5, 命题 4），断言由此得证。
\item
  进而假设 Z 在每点的余维数至少为 3；于是命题 6 及其推论的假设满足，因为沿用前段的记号，E-F 单连通（I, p. 128, 例 4）。
\end{enumerate}

由于微分流形是可解圈空间（IV, p. 347），本节的结果有如下的特例。

设 Y 是一个局部有限维微分流形，Z 是 Y 的一个闭子流形，i 是 Y−Z 到 Y 中的典范单射。

\begin{enumerate}
\def\labelenumi{\alph{enumi})}
\item
  若 Z 在 Y 中的余维数在每点 ≥ 1，则流形 Y−Z 在 Y 中稠密，且映射 \(\pi_{0}(i)\) 是满射的。
\item
  若 Z 在 Y 中的余维数在每点 ≥ 2，则映射 \(\pi_{0}(i)\) 是双射的，且对 Y-Z 的每一点 x，映射 \(\pi_{1}(i,x)\) 是满射的。
\item
  若 Z 在 Y 中的余维数在每点 \(\geqslant 3\)，则对 Y-Z 的每一点 x，映射 \(\pi_{0}(i)\) 与 \(\pi_{1}(i,x)\) 都是双射的。
\end{enumerate}

